\documentclass[11pt]{article}
\usepackage[margin=1in]{geometry}
\usepackage[T1]{fontenc}
\usepackage{lmodern}
\usepackage{microtype}
\usepackage{amsmath}
\usepackage{amssymb}
\usepackage{hyperref}
\ExplSyntaxOn
\ior_new:N \g_binaml_metadata_ior
\tl_new:N \g_binaml_title_tl
\seq_new:N \g_binaml_authors_seq
\seq_new:N \g_binaml_author_emails_seq
\tl_new:N \g_binaml_date_tl
\tl_new:N \g_binaml_status_tl
\cs_new_protected:Npn \binaml_read_metadata:
  {
    \ior_open:Nn \g_binaml_metadata_ior { metadata.yaml }
    \ior_str_map_inline:Nn \g_binaml_metadata_ior
      {
        \regex_match:nnTF { \A title: } { ##1 }
          {
            \tl_gset:Nn \g_binaml_title_tl { ##1 }
            \regex_replace_once:nnN
              { \A title: \s* " } { } \g_binaml_title_tl
            \regex_replace_once:nnN
              { " \s* \Z } { } \g_binaml_title_tl
          }
          {
            \regex_match:nnTF { \A \s* email: } { ##1 }
              {
                \tl_set:Nn \l_tmpa_tl { ##1 }
                \regex_replace_once:nnN
                  { \A \s* email: \s* " } { } \l_tmpa_tl
                \regex_replace_once:nnN
                  { " \s* \Z } { } \l_tmpa_tl
                \seq_gput_right:NV \g_binaml_author_emails_seq \l_tmpa_tl
              }
              {
                \regex_match:nnTF { \A date: } { ##1 }
                  {
                    \tl_gset:Nn \g_binaml_date_tl { ##1 }
                    \regex_replace_once:nnN
                      { \A date: \s* } { } \g_binaml_date_tl
                  }
                  {
                    \regex_match:nnTF { \A status: } { ##1 }
                      {
                        \tl_gset:Nn \g_binaml_status_tl { ##1 }
                        \regex_replace_once:nnN
                          { \A status: \s* } { } \g_binaml_status_tl
                      }
                      {
                        \regex_match:nnT { \A \s* - \s* name: } { ##1 }
                          {
                            \tl_set:Nn \l_tmpa_tl { ##1 }
                            \regex_replace_once:nnN
                              { \A \s* - \s* name: \s* } { } \l_tmpa_tl
                            \seq_gput_right:NV \g_binaml_authors_seq \l_tmpa_tl
                          }
                      }
                  }
              }
          }
      }
    \ior_close:N \g_binaml_metadata_ior
  }
\binaml_read_metadata:
\NewDocumentCommand \BinamlTitle { } { \tl_use:N \g_binaml_title_tl }
\NewDocumentCommand \BinamlAuthors { }
  {
    \seq_map_indexed_inline:Nn \g_binaml_authors_seq
      {
        ##2 \\
        \href
          {mailto:\seq_item:Nn \g_binaml_author_emails_seq { ##1 }}
          {\seq_item:Nn \g_binaml_author_emails_seq { ##1 }}
        \int_compare:nNnT { ##1 } < { \seq_count:N \g_binaml_authors_seq }
          { \and }
      }
  }
\NewDocumentCommand \BinamlDate { }
  {
    \tl_use:N \g_binaml_date_tl
    \str_if_eq:eeT { \tl_use:N \g_binaml_status_tl } { draft }
      { \space (Draft) }
  }
\ExplSyntaxOff
\title{\BinamlTitle}
\author{\BinamlAuthors}
\date{\BinamlDate}
\begin{document}
\maketitle
\begin{abstract}
Binaml focuses on continual learning models over streams of binary features
subject to data drift.  Binary features provide a task-agnostic representation
and enable models optimized for memory, latency, and energy efficiency.  This
paper specifies Binaml's current regression model, \texttt{BRegressor}, as an
initial task-specific instantiation.  It adapts online rather than relying on a
stationary train/serve assumption and learns compact compositional
representations from residual feedback.  The model is environment-agnostic: it
does not depend on how the stream is produced.  This paper specifies the
current regression model and reports reproducible prequential comparisons on
versioned synthetic streams.
\end{abstract}

\section{Notation}
\begin{description}
  \item[$d,x_t,y_t$] Input width, binary input, and scalar regression target.
  \item[$\hat y_t,e_t,s_t$] Prediction, residual, and residual-sign label.
  \item[$b,w_r,\eta,\lambda$] Intercept, coefficient of feature id $r$, learning
  rate, and $\ell_2$ decay.
  \item[$r=(\ell,j)$] Feature id with store layer $\ell$ and index $j$.
  \item[$\varphi_r,\psi_r,\mathcal L$] Boolean value, centered value, and live
  feature set for feature id $r$.
  \item[$T\in\{0,\ldots,15\}$] Four-bit truth table of an arity-2 composed
  feature.
  \item[$r_0,r_1$] Ordered input feature ids of a composed feature.
  \item[$B,S,n,K$] Feature-learning batch size, SGD steps per observation,
  feature-batch index, and parent top-$K$ width.
  \item[$C_{\ell},V_{\ell},L$] Live capacity per learned layer, candidate
  capacity per layer, and maximum learned-layer depth.
  \item[$\sigma(r),\tilde\sigma(r)$] Match score and propagated pruning score.
\end{description}

\section{Model}
At time $t$, the model receives $x_t\in\{0,1\}^d$ and later observes a finite
scalar $y_t$.  If useful signal is sparse weighted boolean structure over bits,
a linear head over learned boolean indicators is a direct hypothesis class:
continuous amplitudes, discrete features, and an inspectable composition.
Let $\mathcal L$ be the current set of live feature ids.
The prediction is
\[
\hat y_t
=
b
+
\sum_{r\in\mathcal L}
w_r\,\psi_r(x_t),
\]
where $b\in\mathbb R$, each $w_r\in\mathbb R$, and each
$\varphi_r(x_t)\in\{0,1\}$ is centered before entering the linear head:
\[
\psi_r(x_t)=2\varphi_r(x_t)-1\in\{-1,1\}.
\]

\paragraph{Layer zero.}
Input coordinates have feature ids
$r=(0,j)$ for $j\in\{0,\ldots,d-1\}$, with
\[
\varphi_{(0,j)}(x)=x_j.
\]
Their initial coefficients are zero.

\paragraph{Composed features.}
Composition starts from the most basic nontrivial boolean learning unit:
two-input tables.  Combining them yields richer features without beginning
from high-arity search.  A composed feature combines two input features
$(r_0,r_1)$ using a four-bit truth table $T\in\{0,\ldots,15\}$.

\section{Feature representation and storage}
Depth is the natural complexity axis for feature composition: it bounds how
far candidates look, keeps evaluation a DAG walk, and makes capacity control
per depth band.  The store layer of a composed feature is
\[
\ell(r)=1+\max\bigl(\ell(r_0),\ell(r_1)\bigr).
\]
Layer zero holds the $d$ source features.  Each learned layer
$\ell\geq 1$ maintains two collections:
\begin{description}
  \item[Live] Optionally occupied slots of composed features.  Indices are
  stable; pruned slots become tombstones and are no longer live feature ids.
  \item[Candidates] Candidate features trained on the current batch and not yet
  admitted to the live set.
\end{description}
A candidate insert requires ordered distinct inputs $r_0<r_1$, valid live or
source parents, and the correct derived layer.  The store permits multiple
features with the same parent pair, allowing truth tables learned on different
batches to coexist.  Newly promoted live coefficients start at zero.  After
promote or prune, coefficients of dead feature ids are dropped and missing
live feature ids receive weight zero.

\subsection{Runtime layout}
The implementation separates stable feature identity from compact evaluation.
Each id is the pair \texttt{FeatureId(layer, index)}.  Source scores occupy one
contiguous vector.  Learned layers occupy a vector of layer records; each
record has a live vector of optional feature records and a candidate vector.
An absent live entry is a tombstone, so pruning does not renumber an id that a
later feature may reference.  A feature record stores two parent ids, one
four-bit \texttt{u8} truth table, and one \texttt{u8} match score.  The eight
\texttt{u8} counters used to learn a table cover every
$(u,v,s)\in\{0,1\}^3$ combination, which limits the feature-learning batch
size to 255.

Before evaluation, the live store is flattened in layer order into a
\texttt{FeatureLayout}: one vector of feature ids and one same-length vector of
nodes.  A source node stores its input coordinate; a composed node stores the
two already evaluated node positions and its truth table.  Thus each input is
evaluated in one forward topological pass.  During \texttt{observe}, the current
sample's activation row is computed once and reused for $S$ linear updates.
Feature learning accumulates residual signs over non-overlapping batches of $B$
binary inputs.  Coefficients are kept separately in a map from feature id to
\texttt{f64}; after each structural update, tombstoned weights are removed and
new live ids receive weight zero.

\begin{center}
\scriptsize
\begin{tabular}{@{}p{0.20\textwidth}p{0.47\textwidth}p{0.25\textwidth}@{}}
\hline
\textbf{Type} & \textbf{Fields or representation} & \textbf{Role} \\
\hline
\texttt{FeatureId} & \texttt{\{layer: usize, index: usize\}} & stable source or learned-feature identity \\
\texttt{Feature} & \texttt{\{inputs: [FeatureId; 2], truth\_table: u8, score: u8\}} & composed boolean feature and pruning score \\
\texttt{FeatureLayer} & \texttt{\{live: Vec<Option<Feature>>, candidates: Vec<Feature>\}} & one learned depth band; empty live slots are tombstones \\
\texttt{FeatureStore} & \texttt{\{source\_scores: Vec<u8>, layers: Vec<FeatureLayer>\}} & owns all source, live, and candidate features \\
\texttt{FeatureLearningConfig} & \texttt{\{batch\_size, parent\_top\_k, features\_per\_layer, candidate\_capacity, max\_layers\}} & feature-search limits and parent-selection width \\
\texttt{FeatureLearner} & \texttt{\{config, store, has\_previous\_batch: bool\}} & scores, promotes, prunes, and generates features \\
\texttt{FeatureLayout} & \texttt{\{references: Vec<FeatureId>, nodes: Vec<FeatureNode>\}} & compact layer-ordered evaluation plan \\
\texttt{FeatureNode} & \texttt{Source(usize)} or \texttt{Composed\{first: usize, second: usize, truth\_table: u8\}} & source coordinate or compiled composed operation \\
\texttt{BRegressor} & learner, coefficient \texttt{HashMap<FeatureId, f64>}, intercept \texttt{f64}, feature-batch buffers, layout, count & online model state \\
\hline
\end{tabular}
\end{center}

\section{Online linear-head updates}
Composed-feature identity is discrete; coefficients are continuous, so the two
are updated separately.  For each observation $(x_t,y_t)$, evaluate the current
live features once, compute the residual sign with the current coefficients,
then perform $S$ single-sample squared-error gradient steps on that activation
row.  For one such step with pre-step residual $e_t=y_t-\hat y_t$,
learning rate $\eta>0$, and decay $\lambda\geq 0$, apply
\begin{align*}
w_r &\leftarrow (1-\eta\lambda)\,w_r
&&\text{for all live }r,\\
b &\leftarrow b+\eta e_t,\\
w_r &\leftarrow w_r+\eta e_t\psi_r(x_t).
\end{align*}
There is no gradient through boolean structure: SGD tracks residual amplitude
at the sample level, while structure search runs when enough residual-sign
evidence exists.  The residual sign
\[
s_t=\mathbf{1}\{e_t\geq 0\}
\]
is computed with the current coefficients before the associated linear updates
and appended to the current feature-learning batch.  After every $B$ new
observations, that batch is consumed for feature learning and cleared.

\section{Residual-sign supervision and batching}
Feature learning targets residual signs, not $y_t$ itself.  Residual signs are
a binary supervision bit for the discrete learner, while the linear head
absorbs magnitude---simple parts combined, not one joint high-complexity
update.  Exact truth-table counting needs a bounded window; batching amortizes
candidate search and gives a natural train/hold-out split across successive
windows.  Batch $n$ is
\[
\mathcal B_n
=
\bigl\{(x_{t},s_{t})\bigr\}_{t\in I_n},
\qquad |I_n|=B,
\]
where $I_n$ is the contiguous block of $B$ observations since the previous
structural update.
Processing $\mathcal B_n$ has three phases when a previous batch exists:
\begin{enumerate}
  \item Score every source, live feature, and candidate on
  $\mathcal B_n$ by match count against residual signs.
  \item Promote every candidate into live, then prune each learned
  layer to capacity $C_{\ell}$.
  \item Learn new candidates on $\mathcal B_n$.
\end{enumerate}
On the first batch only the candidate-learning phase runs.  Candidates trained
on $\mathcal B_n$ therefore receive their first scores on $\mathcal B_{n+1}$
before promotion: a one-batch hold-out across successive residual-sign
windows.

\paragraph{Match score.}
For a boolean stream $z_{1:B}$ and residual signs $s_{1:B}$,
\[
\sigma=\sum_{i=1}^{B}\mathbf{1}\{z_i=s_i\}.
\]
Sources use their coordinate columns; composed features use their evaluated
bitstreams.

\section{Truth-table learning}
Given two parent bitstreams $u_{1:B}$, $v_{1:B}$ and residual signs
$s_{1:B}$, the learner returns the Hamming-error-minimizing four-bit table.
Its state is a fixed array of eight \texttt{u8} counters, one for each
combination of the two parent values and residual sign.  The learned result is
a four-bit truth table represented by one \texttt{u8}.  The counters record
the frequencies needed to select each table output by majority, with ties
resolving to zero.  As one counter may receive every observation, the
\texttt{u8} counters bound the admissible batch size to $B\leq255$.

\section{Candidate feature generation}
For each layer $\ell=1,\ldots,L$, while candidate capacity $V_{\ell}$ remains:
\begin{enumerate}
  \item Let $P$ be the top-$K$ live feature ids in layer $\ell-1$ by match
  score, breaking ties by feature-id order.
  \item Form every distinct pair $\{r_a,r_b\}\subset P$, written with
  $r_0<r_1$.
  \item For each pair, learn a truth table on the current batch and insert it
  into candidates with initial score zero, stopping if candidate capacity is
  reached.  A pair may produce a new candidate even when that pair already
  appears in the layer.
\end{enumerate}
If $P$ has fewer than two feature ids, layer $\ell$ is skipped.
Layer depth never exceeds $L$.

\section{Feature promotion and pruning}
After scoring on the following batch, every candidate at each learned layer is
appended to that layer's live list.  A parent can look weak alone yet be
required by useful children, so pruning uses max-with-descendants scores.
If a layer then exceeds capacity $C_{\ell}$, live feature ids are sorted by
ascending propagated score
\[
\tilde\sigma(r)
=
\max\!\Bigl(
\sigma(r),\;
\max_{r'\in\mathrm{children}(r)}\tilde\sigma(r')
\Bigr),
\]
where $\mathrm{children}(r)$ are live features that take $r$ as an input, and
ties break by feature-id order.  The lowest
$|\mathrm{live}|-C_{\ell}$ feature ids are tombstoned.  Tombstoning cascades
to all dependents, so no live feature retains a dead parent.

\section{Online learning procedure}
\noindent\textbf{Configuration}
\begin{description}
  \item[\textbf{Linear}] Learning rate $\eta$, decay $\lambda$, and $S$ SGD
  steps per observation.
  \item[\textbf{Batch structure}] Feature batch size $B$, parent top-$K$ width
  $K$, per-layer live capacity $C_{\ell}$, candidate capacity $V_{\ell}$,
  and maximum depth $L$.
\end{description}

\noindent\textbf{Initialize:} empty learned layers, $b\leftarrow 0$,
$w_{(0,j)}\leftarrow 0$ for all source indices, and empty feature-batch
buffers.

\noindent\textbf{For} each observation $(x_t,y_t)$:
\begin{enumerate}
  \item Validate the input width and finite target.
  \item Evaluate live $\varphi_r(x_t)$ through the current layout and predict
  $\hat y_t=b+\sum_{r\in\mathcal L}w_r\psi_r(x_t)$.
  \item Compute $s_t=\mathbf{1}\{y_t-\hat y_t\geq 0\}$ and append $(x_t,s_t)$
  to the current feature batch.
  \item Apply $S$ single-sample linear updates using the current activation
  row.
  \item After $B$ accumulated observations, score the previous candidates,
  promote and prune live features, generate candidates from the current
  residual-sign batch, synchronize coefficients, rebuild the layout, and clear
  the feature batch.
\end{enumerate}

\paragraph{Cost.}
Prediction and activation evaluation are linear in the number of live source
and composed features.  Each linear step costs one live-feature count.
Structural work occurs once per feature batch: candidate table learning scans
each selected parent pair over $B$ rows, while pruning recursively scans later
layers to preserve a useful descendant's parents.  The design trades occasional
tombstones for stable ids and bounded live and candidate capacity per layer.

\section{API and integration boundaries}
The \texttt{binaml.models} package owns the online model and has no
environment dependency: the goal is a generic model, decoupled from where the
data comes from and how it was produced.  The public surface is
\texttt{BRegressor}, which implements
\texttt{predict(features)} and \texttt{observe(features, target)} for binary
feature vectors of width $d$.  The prediction call may retain pairing state;
\texttt{observe} must follow a preceding \texttt{predict}.  Feature learning
is internal to \texttt{observe} once a batch fills.  Evaluation packages
compose models with environments through predict-then-observe prequential
evaluation.  Named benchmarks only compose these independent components.

\section{Prequential evaluation protocol}
For each emitted $(x_t,y_t)$, an evaluator first records the model
prediction, then exposes the target through \texttt{observe}.  Mean squared
error is computed from the recorded pre-update predictions.  Structure
updates occur only after a full feature batch fills inside \texttt{observe};
residual signs are measured with the current coefficients before the
corresponding linear updates.
Hyperparameters, implementation version, and the online protocol are required
to interpret any reported trajectory.

\section{Evaluation}
The benchmark inputs and model configuration are versioned in this paper's
benchmark directory.  The task uses 1,000 observations and seeds
$\{0,1,2,3,4\}$ with 32 input bits,
12 fixed target components, noise standard deviation $0.1$, and the same
function, DAG, and target-scale settings.  Input, gate, and intercept
distributions each resample independently with probability $0.01$.

All models receive the scenario width.  The feature regressor uses learning
rate $\eta=10^{-3}$, $\ell_2$ decay $10^{-4}$, feature batch size $B=32$,
three single-sample SGD steps, parent top-$K=8$, 32 live features per learned
layer, candidate capacity 32, and depth two.  Every source and composed
activation enters the linear head as $2\varphi_r(x)-1$, and candidate pairs are
distinct members of the top-$K$ previous layer only.  The linear baseline is a
Python+NumPy replay-batch implementation using learning rate $0.03$, the same
decay, replay size 32, and step count, with uncentered binary inputs.  The MLP
baseline is a Scikit-learn wrapper using replay size 32 and three
\texttt{partial\_fit} passes; one 50-unit ReLU hidden layer; scikit-learn SGD
with constant learning rate $0.003$, $\alpha=10^{-4}$,
\texttt{power\_t}=0.5, shuffling enabled, random state zero, tolerance
$10^{-4}$, zero momentum, no Nesterov momentum, and no early stopping.  Its
remaining configured values are validation fraction $0.1$, $\beta_1=0.9$,
$\beta_2=0.999$, $\epsilon=10^{-8}$, 10 no-change iterations, and disabled
verbose output; the wrapper fixes one iteration and disables warm start.

Tables report the mean $\pm$ standard error over five seeds.  MSE is the
optimized loss; RMSE expresses the same error in target units.  Total time is
the prequential predict-plus-observe time for a 1,000-sample trajectory.

\begin{center}
\footnotesize
\begin{tabular}{lrrr}
\hline
\textbf{Model} & MSE & RMSE & total s \\
\hline
BRegressor (ours) & $\boldsymbol{0.494\pm0.022}$ & $\boldsymbol{0.702\pm0.015}$ & $\boldsymbol{0.274\pm0.053}$ \\
Linear SGD (Python+NumPy) & $0.745\pm0.015$ & $0.863\pm0.008$ & $0.172\pm0.005$ \\
MLP SGD (Scikit-learn) & $0.842\pm0.024$ & $0.917\pm0.013$ & $3.133\pm0.201$ \\
\hline
\end{tabular}
\end{center}

\section{Limitations and scope}
This document specifies the feature regressor used by Binaml.  It does not
claim optimality of residual-sign learning.  The reported comparison is
limited to the versioned synthetic task, fixed configurations, and
reproducible baseline runs described above.
\end{document}
